\section{Higher risk aversion: the frontier for \texorpdfstring{$\mu\in\{3,5\}$}{mu = 3, 5}}
\label{sec:beyond}

Everything in Section~\ref{sec:arch} except Theorem 1 goes through for CRRA utility with any
$\mu>0$: the corner, the minimal MPC with $\kappa=1-\Thorn/R$, the cap slack, the Doeblin
uniqueness and convergence, the coupling, the homogeneity, single crossing. The formal statement is
\lean{crra\_equilibriumRate\_unique\_of\_marginal}, and it carries exactly one hypothesis that is
not discharged: Theorem 1. This section says what that hypothesis is, why it fails at Aiyagari's
numbers, what it becomes when the proof is tightened, and where that leaves the problem.

\subsection{Theorem 1 reduces to an elasticity condition}\label{sec:elasticity}

\begin{proposition}[Saving rises with the rate under an elasticity condition]\label{prop:elas}
Let $r_1\le r_2$ be admissible, with positive consumption and a slack cap at both rates, and let
$u'$ be strictly decreasing. Fix $(a,z)$ with $a\ge0$ and let $b=g_1(a,z)$. If
\begin{equation}\label{eq:EC}
  R_1\,u'\bigl(c_1(b,z')\bigr)\ \le\ R_2\,u'\bigl(c_2(b,z')\bigr)\qquad\text{for every }z',
\end{equation}
then $g_1(a,z)\le g_2(a,z)$.
\formal{IncomeFluctuation.policy\_le\_policy\_withRate\_of\_marginal\_at}
\end{proposition}

The proof is four lines of Euler inequalities and holds for any period utility. Suppose
$g_2(a,z)<g_1(a,z)=b$. Then the higher-rate household has room under the cap, so
$\beta R_2\,\E[u'(c_2(g_2,\cdot))]\le u'(c_2(a,z))$; the lower-rate household has room above the
floor, so $u'(c_1(a,z))\le\beta R_1\,\E[u'(c_1(b,\cdot))]$; it consumes strictly less today, so
$u'(c_2(a,z))<u'(c_1(a,z))$; and consumption rises with assets, so
$u'(c_2(b,\cdot))\le u'(c_2(g_2,\cdot))$. Chaining, $R_2\E[u'(c_2(b,\cdot))]<R_1\E[u'(c_1(b,\cdot))]$,
which \eqref{eq:EC} forbids. For CRRA, \eqref{eq:EC} is
\[
  \Bigl(\frac{c_2(b,z')}{c_1(b,z')}\Bigr)^{\mu}\le\frac{R_2}{R_1},
\]
an elasticity of consumption to the gross rate of at most $1/\mu$ (\lean{crra\_elasticity\_iff}).
It says that the income effect of the rate on consumption does not outweigh the human-wealth
effect. Uniqueness on $[r_{lo},r_{hi}]$ follows from \eqref{eq:EC} at every state and every pair of
rates (\lean{crra\_equilibriumRate\_unique\_of\_marginal}), and, since the coupling only visits
states the economy visits, from \eqref{eq:EC} on any interval $[0,\bar a']$ that the lower-rate
policy maps into itself (\lean{crra\_equilibriumRate\_unique\_of\_marginal\_on},
Section~\ref{sec:invariant}).

\subsection{The condition fails for the rich when \texorpdfstring{$\mu>1$}{mu > 1}}

By \citet{MaToda2022}, the consumption function is asymptotically linear, $c(x)/x\to\kappa(R)$,
with $\kappa(R)=1-(\beta R)^{1/\mu}/R$ the minimal MPC. For $\mu>1$, $\kappa$ rises with $R$. So at
high wealth $c_2(b)/c_1(b)\to\kappa_2R_2/(\kappa_1R_1)$ and
\[
  \Bigl(\frac{\kappa_2R_2}{\kappa_1R_1}\Bigr)^{\mu}
  =\Bigl(\frac{\kappa_2}{\kappa_1}\Bigr)^{\mu}\Bigl(\frac{R_2}{R_1}\Bigr)^{\mu}>\frac{R_2}{R_1},
\]
so \eqref{eq:EC} fails for all large $b$ (\lean{crra\_elasticity\_fails\_of\_asymptotic\_mpc}).
Saving of the rich nonetheless rises with the rate, since $g\approx\Thorn b$ with
$\Thorn=(\beta R)^{1/\mu}$ increasing in $R$; the sufficient condition does not see it. With the
linear approximation $c\approx\kappa(R)(Rb+H)$ and $H\approx\bar Y/r$, the human-wealth effect
dominates and consumption falls with the rate for $b$ below about $13\bar Y$ at $\mu=3$ and
$7\bar Y$ at $\mu=5$, where \eqref{eq:EC} holds; above, it fails. The stationary distribution
reaches far beyond those levels, because the top income type keeps accumulating until wealth of
order $y_{\max}/(1-\Thorn)$, so the invariant-interval form does not rescue the argument at
Aiyagari's numbers.

\subsection{The chain throws a term away}\label{sec:slack}

The supposition $g_2(a,z)<g_1(a,z)$ gives not $c_2(a,z)>c_1(a,z)$ but
\[
  c_2(a,z)>c_1(a,z)+(r_2-r_1)\,a,
\]
and the proof above uses only the weaker inequality. Keeping the extra interest income, the chain
closes under a weaker condition.

\begin{proposition}[The slack elasticity condition]\label{prop:slack}
In the setting of Proposition~\ref{prop:elas}, with $u'>0$, if
\begin{equation}\label{eq:SEC}
  R_1\,u'\bigl(c_1(b,z')\bigr)\,u'\bigl(c_1(a,z)+(r_2-r_1)a\bigr)
  \ \le\ R_2\,u'\bigl(c_2(b,z')\bigr)\,u'\bigl(c_1(a,z)\bigr)\qquad\text{for every }z',
\end{equation}
then $g_1(a,z)\le g_2(a,z)$. For CRRA, \eqref{eq:SEC} reads
\[
  \Bigl(\frac{c_2(b,z')}{c_1(b,z')}\Bigr)^{\mu}\le\frac{R_2}{R_1}
  \Bigl(1+\frac{(r_2-r_1)\,a}{c_1(a,z)}\Bigr)^{\mu}.
\]
\formal{IncomeFluctuation.policy\_le\_policy\_withRate\_of\_slack\_at; crra\_slack\_elasticity\_iff}
\end{proposition}

The new factor is at least one and grows with $a/c_1(a)$, which is largest exactly where
\eqref{eq:EC} fails. In the high-wealth limit, with $L_i=\kappa_iR_i=R_i-\Thorn_i$ the asymptotic
slopes and $b\approx\Thorn_1a$, $c_1(a)\approx L_1a$, both sides have limits and the inequality
between them is exact:

\begin{proposition}[The slack condition holds in the homogeneous limit for every $\mu$]
\label{prop:homog}
For $0<R_1\le R_2$, $\Thorn_1\le\Thorn_2$, $\Thorn_1<R_1$, $\Thorn_2<R_2$ and $\mu\ge0$,
\[
  \Bigl(\frac{R_2-\Thorn_2}{R_1-\Thorn_1}\Bigr)^{\mu}\ \le\
  \frac{R_2}{R_1}\Bigl(\frac{R_2-\Thorn_1}{R_1-\Thorn_1}\Bigr)^{\mu},
\]
with strict inequality when $\Thorn_1<\Thorn_2$ and $\mu>0$. Consequently, along tomorrow's
assets, if $c_2(b(a))/c_1(b(a))\to L_2/L_1$ and $c_1(a)/a\to L_1$, the CRRA slack condition holds
for all large $a$.
\formal{crra\_slack\_elasticity\_of\_patience; crra\_slack\_elasticity\_holds\_of\_asymptotic\_mpc}
\end{proposition}

The proof is one line: $R_2-\Thorn_2\le R_2-\Thorn_1$ because $\Thorn_1\le\Thorn_2$, and
$R_2/R_1\ge1$ is spare. In logarithmic terms per unit of rate difference the two sides grow at
rates $\mu(\kappa'R+\kappa)/(\kappa R)$ and $1/R+\mu/(\kappa R)$, whose difference is exactly
$1/(\kappa R)$ for every $\mu$: about $25$ at Aiyagari's numbers. Table~\ref{tab:slack} gives the
values.

\begin{table}[t]
\centering
\caption{The slack elasticity condition at Aiyagari's numbers, $r\approx4\%$. Rows two and three
are the logarithmic growth rates of the two sides of \eqref{eq:SEC} per unit of rate difference in
the high-wealth limit; the condition holds when the left is below the right. Row four is the wealth
below which consumption falls with the rate, where the unweakened condition \eqref{eq:EC} already
holds.}
\label{tab:slack}
\begin{tabular}{lcc}
\toprule
 & $\mu=3$ & $\mu=5$\\
\midrule
minimal MPC $\kappa$ & $0.0390$ & $0.0388$\\
left side, deep tail & $50$ & $100$\\
right side, deep tail & $75$ & $125$\\
wealth below which \eqref{eq:EC} holds (multiples of mean income) & $\approx13$ & $\approx7$\\
\bottomrule
\end{tabular}
\end{table}

So at every wealth level the slack condition is, numerically, true: in the body because
consumption falls with the rate, in the tail by Proposition~\ref{prop:homog} with margin, and in
between with margin as well. The reduction theorems with the slack condition in place of
\eqref{eq:EC} are \lean{equilibriumRate\_unique\_of\_slack},
\lean{equilibriumRate\_unique\_of\_slack\_on} and
\lean{crra\_equilibriumRate\_unique\_of\_slack\_on}.

\subsection{What remains}\label{sec:frontier}

To prove \eqref{eq:SEC} from primitives one needs the ratio $c_2(b)/c_1(b)$ of consumption at two
rates to a relative precision of the order of $r_2-r_1$, uniformly in $b$ on the support, since
that is the size of the margin. Two available routes do not deliver it. The sandwich
$\kappa x\le c(x)\le\kappa x+C$, with $C$ of the order of mean income, has a gap of order one and
so pins the ratio only for rate pairs far enough apart, whereas uniqueness must handle all pairs.
The fixed-point perturbation bound for the Coleman operator in the marginal-utility metric of
\citet{LiStachurski2014}, which we formalised (Section~\ref{sec:aux}), gives a bound that is
uniform in wealth in the wrong sense: it loses a factor of the inverse contraction modulus,
$1/(1-\Thorn)\approx2000$ at these numbers, and is vacuous where $u'$ is small. What would be
needed is a Lipschitz estimate in the rate for the consumption function that is relative rather than
absolute, that is, comparative statics of $c$ in $r$ uniform on the support. We know of no such
result. The status of $\mu\in\{3,5\}$ is therefore: reduced to an inequality on two consumption
functions that is numerically true everywhere, with an explicit margin, and unproved.

Two further reductions are recorded for completeness. The rescaling proof of Section~\ref{sec:thm1}
uses relative risk aversion at most one only through a range condition on the true consumption
function of the higher-rate economy, $(1-c_2(0)/c_2(ta))(\rho-1)\le\rho^{1/\mu}-1$ with
$\rho=R_2/R_1$, roughly $c(a)\lesssim\frac{\mu}{\mu-1}c(0)$ (\lean{policy\_mono\_withRate\_crra\_of\_range}).
It is weaker than $\mu\le1$ but fails at moderate wealth. And Theorem 1 holds at equal cash on
hand for all $a$: $g_1(a,z)\le g_2(ta,z)$ with $t=R_1/R_2$ (\lean{policy\_le\_policy\_scaled\_withRate\_of\_marginal\_at}).

\subsection{The invariant-interval form}\label{sec:invariant}

Since Theorem 1 uses the condition only at tomorrow's assets and the coupling only along coupled
paths, if $[0,\bar a']$ is invariant under the lower-rate policy then the condition on $[0,\bar a']$
alone orders the policies there, the coupling started at zero assets never leaves the interval, and
uniqueness follows (\lean{AboveSet}, \lean{push\_aboveSet\_eq\_zero},
\lean{equilibriumRate\_unique\_of\_marginal\_on}). This is the honest form of the reduction: the
condition is needed on the ergodic support and nowhere else. At Aiyagari's numbers the support is
too long for the unweakened condition, as noted; for the slack condition the invariant form is
available but, by Proposition~\ref{prop:homog}, not needed.
